\section{A countably infinite Hilbert space with quantitative names}
\label{sec:infinitehilbert}

Fix the discrete basis $e_0,e_1,\ldots$. A finite vector is a finite list of
Gaussian rationals, with trailing zeroes ignored. Write
$\delta_k=2^{-k}$, $P_N$ for restriction to $0\le i<N$, and
$\lambda(C)=\min\{h\le C:2^h\ge C\}$ for positive integral $C$.
Ceil a positive rational before applying $\lambda$. Squared norms of finite
vectors are rational. Thus even a test $\|u-v\|\le r$ for nonnegative
rational $r$ is a rational test. Inner products are conjugate-linear in the
first slot. All algorithms below apply to fixed PR presentations or uniformly
to families carrying PR algorithms, in the schematic sense of iPRA.

\subsection{Vectors, tails and completeness}

A name consists of an integer $M\ge1$ and a PR function $s\mapsto x_s$ of
finite vectors with
\[
 \|x_s\|\le M,\qquad \|x_s-x_t\|\le\delta_s\quad(t\ge s).
\]
Equality of two names means $\|x_s-y_s\|\le 2\delta_s$ for every $s$.
This relation is transitive: compare both names at an arbitrarily late
stage with an intermediate name, then return to stage $s$. The resulting
bound is $2\delta_s+4\delta_t$ for every $t\ge s$; a strict rational
violation of $2\delta_s$ supplies a finite $t$ excluding it. Denote the
setoid by $\ell^2_{\rm PR}$. An estimate such as $\|x-x_s\|\le\delta_s$
means the corresponding estimates at finite stages with their rounding
allowances. In computations one can target error $\delta_{k+2}$ at stage $k$:
two such approximations to the same vector differ by at most $\delta_k$
at stages $k$ and $t\ge k$.

The finite Cauchy--Schwarz inequality follows, including the zero-vector
case, from the polynomial identity
\[
 \|a\|^2\|b\|^2-|\langle a,b\rangle|^2
   =\sum_{i<j}|a_i b_j-a_j b_i|^2.
\]
It gives the triangle inequality by expanding a square. Addition,
multiplication by a named scalar, the norm, and the inner product pass to
names. For example, when inputs have bounds $M,N$, the inner-product
error made by using both stage-$s$ vectors is at most $(M+N)\delta_s$;
the norm error is at most $\delta_s$ before rational square-root rounding.
For scalar products use
$\|a x-b y\|\le |a|\|x-y\|+|a-b|\|y\|$. A stage
$k+3+\lambda(1+M+N)$ therefore suffices for the inner product and
for the analogous scalar-vector product with these envelopes. Finite norm
rounding uses bisection of a nonnegative rational square root on $[0,M]$.
The finite identities extend by these bounds, including positivity,
conjugate symmetry and the parallelogram law. If $\|x\|=0$ then
$\|x_s\|\le\delta_s$ and $x=0$.

\begin{proposition}[Two presentations of the basis expansion]
A vector name gives uniformly named coordinates and a PR cutoff $L(k)$
such that every finite section of its tail beyond $L(k)$ has norm at
most $\delta_k$. Conversely, uniformly named coordinates with this
finite-tail certificate and a fixed bound $M$ on the norms of all finite
initial sections give a vector name.
\end{proposition}
\begin{proof}
Read a coordinate from $x_s$ with error at most $\delta_s$. Take $L(k)$
equal to the length of $x_k$. Every finite selection of entries beyond it
is a contraction of $x-x_k$. Conversely take $N=L(k+3)$ and approximate the
first $N$ coordinates to complex absolute error
$\delta_{k+3+\lambda(1+N)}$ each. The resulting vector $q_k$ has error at
most $\delta_{k+3}$ in that section, since the norm is at most the sum
of coordinate errors. For $s\ge k$, compare $q_k,q_s$ inside the longer
of their two sections; each incurs its rounding error and at most its
certified tail error. Their difference has norm at most $4\delta_{k+3}$.
The integer $M+1$ bounds their norms. Coordinate recovery follows by
comparing in a section containing the desired index, using the tail estimate
even when that index lies beyond the chosen cutoff. Reconstruction of a
given name is equal to the original by the same estimates. A monotone
cutoff may always be obtained by taking finite running maxima.
\end{proof}

The same estimates give the relevant completeness principle. Suppose
$y^{(j)}$ is a uniformly named family, and a PR function $m$ certifies
$\|y^{(a)}-y^{(b)}\|\le\delta_r$ for $a,b\ge m(r)$.
Replace $m$ by running maxima. For stage $k$ take stage $k+3$ of
$y^{(m(k+3))}$. Two such samples differ by at most $3\delta_{k+3}$.
Their norms are bounded, for example, by the bound of $y^{(m(0))}$
plus $2$. Comparison against a later sample proves convergence with the
given modulus. This diagonal construction uses a supplied joint PR sampler.
A collection of individually PR names need not have one.

The tail certificate carries mathematical information. Let $h(e,j)$ be the
decidable predicate that machine $e$ first halts at step $j$. The sequence
$a_j=h(e,j)$ has at most one nonzero coordinate, and every finite partial
sum of squared moduli is at most one. A uniform computable cutoff $L_e(1)$
would decide halting by simulating up to that cutoff: a later $1$ would give
a finite tail of norm one. This external reduction also excludes a PR rule
producing such cutoffs from the coordinate algorithm and norm upper bound.
It explains precisely why those two pieces of data alone do not give the
above Hilbert-space presentation.

\subsection{Bounded maps and tensors}

A bounded-map presentation for $A$ gives uniformly named columns
$c_j=Ae_j$ and an integer $B\ge1$ with the certificate
\[
 \Big\|\sum_j q_jc_j\Big\|\le B\|q\|
 \qquad\text{for every finite Gaussian-rational vector }q .
 \tag{H1}\label{eq:H1}
\]
The inequality expands to universal precision tests. For an actual
adjointable map one supplies a second bounded-map presentation with columns
$d_i$, and the identities $(c_j)_i=\overline{(d_i)_j}$.
When needed enlarge $B$ to cover both bounds. There is no operation
of searching for a norm or an adjoint in these definitions.

\begin{theorem}[Quantitative extension and adjoint]
These data determine an extensional linear map on vector names, bounded by
$B$. The second presentation acts as its Hilbert-space adjoint.
\end{theorem}
\begin{proof}
For target stage $k$ take $s=k+3+\lambda(B)$, $q=x_s$, and
$C=1+\lceil\sum_j(|\Re q_j|+|\Im q_j|)\rceil$. Return the finite sum
of $q_j$ times the stage-$k+3+\lambda(C)$ columns. Rounding the columns
costs at most $\delta_{k+3}$. Changing from $x_s$ to any subsequent
input stage costs at most $B\delta_s\le\delta_{k+3}$ by
\eqref{eq:H1}. Thus the stage vectors are regular, have bound $BM+1$,
and represent the unique bounded extension, since finite input approximants
are dense with a known rate. Addition and scalar multiplication follow by
comparison on finite inputs and taking the prescribed error to zero.
Pairing the columns on two finite inputs gives
$\langle y,Ax\rangle=\langle A^*y,x\rangle$.
The inner-product error bound extends this identity to names.
In particular, if all columns vanish, the whole map vanishes.
\end{proof}

Columns of $AD$ are $A(De_j)$, the norm bound is the product of the bounds,
and $(AD)^*=D^*A^*$. For a sum take a sum of bounds. An isometry can be
specified by the finite-input identity $\|Aq\|^2=\|q\|^2$; a supplied
unitary additionally has a two-sided inverse, for example by certifying
$A^*A=AA^*=I$ on the columns. These presentations are closed under all
three constructions just given, with nested samplers primitive recursive.

Code a pair of indices using a fixed PR bijection with PR inverses, e.g.\
diagonal enumeration of $\mathbb N^2$ with bounded inversion. The tensor
space consists of regular names of finitely supported arrays in the pair
index. Tensoring finite vectors gives the identities
\[
 \|x\otimes y\|=\|x\|\|y\|,\qquad
 \langle x\otimes y,u\otimes v\rangle
   =\langle x,u\rangle\langle y,v\rangle .
\]
They extend to names because approximating the two factors incurs error
at most $(M+N)\delta_s$ for stage bounds $M,N$. If $A$ is bounded, then
$A\otimes I$ acts on each second-index column of a finite array. Sum the
squares of \eqref{eq:H1} over those columns to obtain bound $B$ on the
array norm. Its columns are $Ae_i\otimes e_j$; the extension theorem
therefore applies. The argument for $I\otimes D$ interchanges the indices.
Composition gives $A\otimes D$ with the product bound; the adjoint and
permutation identities hold on basis arrays and extend. This realizes the
completed separable tensor product with the given bases.

\subsection{Nuclear operators and an intrinsic trace}

For vectors $u,v$ write $[u,v]x=u\langle v,x\rangle$. A nuclear
presentation consists of uniformly named $u_j,v_j$, nonnegative PR
rationals $m_j\ge\|u_j\|\|v_j\|$, an integer $C$ bounding all finite
partial sums of the $m_j$, and a PR tail function $T$ with
\[
  \sum_{a\le j<b}m_j\le\delta_r\qquad(b\ge a\ge T(r)).
 \tag{H2}\label{eq:H2}
\]
It presents $Sx=\sum_j[u_j,v_j]x$. Two presentations represent the same
operator when all matrix entries agree. An upper bound for the classical
trace norm is supplied by $C$; no minimization over decompositions is used.
Finite-rank rational operators, and PR nuclear series equipped with the
displayed certificates, are included.

\begin{theorem}[Trace and bounded observables]
$S$ and $S^*$ have bounded presentations with bounds $\max(1,C)$.
For any bounded-map presentation $A$ with bound $B$, the scalar
\[
  \mathcal T_A(S)=\sum_j \langle v_j,A u_j\rangle
\]
is a regular complex name depending only on the operator represented by
$S$. Thus $\operatorname{Tr}S=\mathcal T_I(S)$ is intrinsic, and
$|\mathcal T_A(S)|\le BC$. If $A$ has a presented adjoint, then
$\operatorname{Tr}(AS)=\mathcal T_A(S)=\operatorname{Tr}(SA)$.
\end{theorem}
\begin{proof}
For $\|x\|\le X$, each term of $Sx$ has norm at most $Xm_j$.
Equation \eqref{eq:H2} provides the Cauchy modulus $T(r+\lambda(X))$
when $X\ge1$. Apply the diagonal construction to finite partial sums;
swap $u_j,v_j$ for $S^*$. Bounds and the adjoint identity follow first
for partial sums and then at every precision.

The scalar tail is at most $B\delta_r$. For an explicit stage-$k$
approximation choose $r=k+3+\lambda(B)$, $J=T(r)$.
Let $U_j,V_j\ge1$ be integer envelopes for $u_j,v_j$.
Take
\[
 h=k+3+\lambda\!\left(1+\sum_{j<J}(BU_j+V_j+1)\right).
\]
Pair finite stage-$h$ samples for $Au_j,v_j$. The first sample has
distance at most $\delta_h$ from a vector of norm at most $BU_j$, hence
the summed inner-product error is at most
$\sum_{j<J}(BU_j+V_j+1)\delta_h$. Regularity follows.

To see independence of the expansion, use a finite basis trace
$D_N=\sum_{i<N}\langle e_i,A S e_i\rangle$. Distributing the finite sum
over finite partial nuclear sums and using their tail bounds gives
$D_N=\sum_j\langle P_Nv_j,A u_j\rangle$ for each fixed $N$.
For $j\ge J$ the total absolute difference from $\mathcal T_A(S)$
is bounded by $B\delta_r$. For $j<J$ take
$h'=k+3+\lambda(1+B\sum_{j<J}U_j)$ and choose $N$ large enough
to contain the supports of the first $J$ samples $(v_j)_{h'}$.
Then $\|(I-P_N)v_j\|\le\delta_{h'}$, and the difference over the
head is at most $B\sum_{j<J}U_j\delta_{h'}$. Both estimates persist
for larger $N$. Two presentations with the same matrix entries have the
same columns, since a named vector with every coordinate zero has norm
zero by its effective tails. Thus they have identical $D_N$; take the
larger of their two cutoffs to prove equality of their traces. Finally
$AS$ has expansion $[Au_j,v_j]$ and $SA$ has expansion
$[u_j,A^*v_j]$. Their majorants are $Bm_j$, with the tail request shifted
by $\lambda(B)$, and the adjoint identity equates their traces.
\end{proof}

\subsection{Density states, discarding a system and observing an outcome}

A convenient positive-state presentation is a unit vector $w$ on a
system and a countable ancillary basis. Finite ancillary bases are included
by padding with zeroes. Write $w_j$ for the system vector obtained by
keeping ancilla index $j$. These slices have joint names by contraction.
Define
\[
  \rho_w=\sum_j[w_j,w_j],\qquad
  \omega_w(A)=\langle w,(A\otimes I)w\rangle .
\]
The presentation specifies the probabilities of experiments on the system.
Distinct amplitudes represent the same mixed state if their density
operators have identical matrix entries.

\begin{proposition}[Density and reduction]
$\rho_w$ has a nuclear presentation, is positive, and has trace one.
For every presented bounded $A$, $\omega_w(A)=\mathcal T_A(\rho_w)$.
Discarding a factor of a multipartite system amounts to moving its
index into the ancilla. This gives a presentation-independent operation
on these states.
\end{proposition}
\begin{proof}
For every finite set of indices, the sum of slice norms squared is at
most $\|w\|^2$ by orthogonal coordinate restriction. Choose a rational
$m_j$ in $[\|w_j\|^2,\|w_j\|^2+\delta_{j+2}]$ by approximating
the square with two extra bits and adding its error. The budget is at
most $1+\sum_j\delta_{j+2}\le 2$. For target $r$ let $J$ exceed
every ancilla index occurring in the finite sample $w_{r+2}$, and
increase it to at least $r+2$. Every finite slice tail beyond $J$
has total squared norm at most $\delta_{r+2}^2$, by contraction
applied to $w-w_{r+2}$. The rounding slacks sum to at most
$2^{-J-1}$. Their sum with $\delta_{r+2}^2$ is below $\delta_r$.
Thus \eqref{eq:H2} holds. Also the finite ancilla projection up to $J$
approximates $w$ in norm to within $\delta_{r+2}$, so summing all
squared slice norms gives exactly one.
For finite projections the pairing formula is just summing over
orthogonal slices. The norm bound for $A\otimes I$ passes it to the
limit. In particular the trace is one and
$\langle x,\rho_w x\rangle=\sum_j|\langle w_j,x\rangle|^2\ge0$.
If an additional system factor is discarded, rearrange the two ancillary
indices with the pairing bijection. For every remaining bounded test $A$
the new expectation equals the old expectation of $A\otimes I$ by
the pairing formula. Testing matrix units recovers every density entry,
which proves independence of the initial purification.
\end{proof}

\begin{theorem}[Every presented density has a purification]
Suppose $S$ has a nuclear presentation, is self-adjoint and positive,
and has trace one. One can construct a uniform unit-vector name $w$ on
system $\otimes$ ancilla such that $\rho_w=S$.
\end{theorem}
\begin{proof}
Take $L=1+C$, so $0\le S\le L^2I$.
The bounded-map polynomial algorithm of Theorem~\ref{thm:q-sqrt}
applied to columns constructs a positive self-adjoint $R$ with $R^2=S$.
For a column target of $\delta_h$ one may take
$m=L^22^{2h+4}$ terms and use the bounded-map extension for the
finite polynomial to a further $\delta_{h+2}$ of error. In detail, $X=I-S/L^2$ is contractive. Successive applications to a
finite vector can each be rounded to $\delta_{h+3}/[L(m+1)]$; the
$j$th iterate then has error at most $j$ times this budget. The absolute
nonconstant coefficients of the square-root polynomial sum to at most
$L$. Thus the variable number of powers uses one PR iteration on finite
vectors and its total error is at most $\delta_{h+3}$.
Set the desired $j$th ancillary slice to $R e_j$. Then
\[
 \sum_{j<N}\|R e_j\|^2
    =\sum_{j<N}\langle e_j,S e_j\rangle=D_N .
\]
The basis-trace proof with $A=I$ gives a PR cutoff $J(r)$ past which
$|D_N-1|\le\delta_r$. The left side is nondecreasing, and its finite
tails from any $N\ge J(r)$ are bounded by $\delta_r$, since each
partial sum is at most its named limit one. At stage $k$ use
$N=J(2k+6)$; assemble approximations to the first $N$ column vectors
of $R$ to error $\delta_{k+3+\lambda(1+N)}$ each.
The omitted slice vector has norm at most $\delta_{k+3}$ in every
finite section, and the finite rounding error is at most $\delta_{k+3}$.
These estimates prove regularity of the assembled name and its norm
equals one by the displayed identity. For a matrix entry of $\rho_w$
the finite slice sum is the corresponding entry of $R P_N R$ by
self-adjointness. On the input $e_i$ its difference from $R^2$ has
norm at most $L\|(I-P_N)R e_i\|$, with PR cutoff from the named
column $R e_i$. Passing to the limit gives $\rho_w=R^2=S$.
\end{proof}

A mixed preparation with PR unit vectors $x_j$ and nonnegative named
probabilities $p_j$ summing to one with PR tail modulus $T$ is included:
use slices $w_j=\sqrt{p_j}\,x_j$. A square root of $p\in[0,1]$ is
obtained by approximating $p$ within $\delta_{2k+6}$, clipping to
$[0,1]$, and bisecting the rational approximation's square root to
$\delta_{k+3}$; $|\sqrt a-\sqrt b|\le\sqrt{|a-b|}$ controls the
first error. Choose $J=T(2k+6)$ and approximate the first $J$ slices
to $\delta_{k+3+\lambda(1+J)}$ each. The assembled error is at most
$\delta_{k+3}$ from the finite approximations and $\delta_{k+3}$ from
the omitted slices. Values $p_j=0$ are allowed. This construction does
not require an eigenbasis for a density operator.

An operation is supplied by a presented isometry $V$ from the input to
output $\otimes$ environment, together with its adjoint when a Heisenberg
map is desired. On a purification apply $V\otimes I$ and move the new
environment index into the ancilla. Its trace is one by isometry. The
output expectation of $A$ is the input expectation of
$V^*(A\otimes I)V$. Consequently the state operation is extensional
and affine, since matrix entries are tested by matrix units. At every
finite matrix level a positive presented test remains positive under
this Heisenberg transformation: for an extra finite factor take the
quadratic form on $(I\otimes V)z$, after permuting the environment
factor. Positivity of a test tensored with the environment identity
follows by summing the nonnegative quadratic forms on finite slices
and applying the explicit tail estimate. The identity test is fixed by
$V^*V=I$. This proves complete positivity and normalization for these
dilations. With only a forward map the pure-vector sampling construction
still works; the adjoint presentation supplies the stated observable
transformation and the argument for extensionality on mixed preparations.

For a countable measurement use a distinguished outcome index in $V$'s
environment. Let $z=(V\otimes I)w$, permuting indices so that the
outcome comes last, and $z_a$ its outcome-$a$ slice. The Born rule and
conditional state become
\[
   p_a=\|z_a\|^2,\qquad
   \rho_{z_a/\sqrt{p_a}}\quad\text{given a witness }p_a\ge\delta_b>0 .
\]
Here $z_a$ retains the output and any unobserved ancillary factors.
The unnormalized slice always exists. For the quotient, $\sqrt{p_a}$
has certified lower bound $\delta_b$, and reciprocation on this region
has Lipschitz bound $2^{2b}$. Rational estimates may first be clipped
to $[\delta_b,1]$. Every PR-decidable event has probability
$\|(I\otimes P_E)z\|^2$, where the projection simply filters the
finite sampled array by its outcome. Orthogonality gives finite
additivity and conditioning on a disjoint union gives the sum of its
unnormalized branch states, as is seen by the slice expansion. The
cutoff determined from the largest outcome in $z_{r+2}$ bounds the
remaining mass by $\delta_{r+2}^2$.
This proves countable normalization and supplies its rate even for
an infinite-outcome device. More generally a jointly PR disjoint family
of events admits effective summation to its union when there is a PR
bound locating the part containing each accepted outcome among the
finitely many retained indices; such a bound is automatic for the
singleton outcome decomposition used here.

\paragraph{Arithmetic content of the estimates.}
Support lengths, maxima of finite lists, bisection depths and all sampling
indices above are PR in their numeric parameters. Finite Pythagorean
identities, rearrangements and norm estimates can be verified as
integer or rational finite-array statements after clearing positive
denominators. Their loops use quantifier-free induction with the free
precision, cutoff and input parameters fixed. Uses of equalities of
completed objects compare arbitrary output precisions and then substitute
the displayed finer precisions. In particular the completion argument
and the trace proof do not perform induction over a quantified
convergence predicate.
